%% Supplementary Information for the sports paper.
\clearpage
\setcounter{figure}{0}
\renewcommand{\thefigure}{S\arabic{figure}}
\setcounter{table}{0}
\renewcommand{\thetable}{S\arabic{table}}

\section*{Supplementary Information}

\noindent This supplement records data provenance, quality-control notes, and the classical derivations of alignment and the $1/(2n)$ cost used in the main text.
Further geometric identities appear in a companion manuscript \parencite{pefstatistics2026}.

\medskip
\noindent\begin{tabular}{@{}p{0.22\linewidth}p{0.34\linewidth}p{0.38\linewidth}@{}}
  \toprule
  \textbf{SI block} & \textbf{Main paper} & \textbf{Contents} \\
  \midrule
  \S1 Data note & Methods & Leagues, seasons, outcome coding \\
  \S2 Quality control & Methods & Bootstrap and leave-one-team-out notes \\
  \S3 Alignment and cost & Theory; Results & Classical derivations \\
  \bottomrule
\end{tabular}

\clearpage

% ---------------------------------------------------------------------------
\refstepcounter{sisection}
\section*{SI Section \thesisection: Data Note}
\label{sec:si_data}

\noindent Rugby KPIs are from the United Rugby Championship ($\PEFrugbyNgames$ matches; $\PEFrugbyNkpis$ KPIs).
Football KPIs are from the English Championship ($\PEFfootNgames$ matches; $\PEFfootNkpis$ KPIs).
Both windows pool seasons 23/24--24/25.
Home is $A$ and away is $B$.
The outcome is home win versus not; draws are coded as not a home win.
Neither league's match-level KPI file is redistributed (\cref{sec:data_availability}).

One row per KPI is written to \path{Alternative Pipeline/pef_numerics/results/E6_sports.csv}.
Second-moment estimates $(\hat\kappa,\hat\rho,\hat\eta)$ and quadrant labels are in \path{scripts/paper_pipeline/outputs/pef_landscape_2season.csv}.
The four pre-specified exemplars are listed in \path{scripts/paper_pipeline/outputs/confirmatory_exemplars.csv}.
Numeric macros in \path{scripts/paper_pipeline/outputs/sports_numbers.tex} are regenerated from those files without refitting.

\clearpage

% ---------------------------------------------------------------------------
\refstepcounter{sisection}
\section*{SI Section \thesisection: Quality Control}
\label{sec:si_qc}

\noindent Team-cluster bootstrap percentile intervals for $(\hat\kappa,\hat\rho,\hat\eta)$ use $B=1{,}000$ home-team resamples.
Away matches of a resampled club can remain in the bootstrap, because the cluster is the home team.
That design matches the landscape pipeline and is separate from leave-one-team-out log-loss (\cref{sec:ml_cv}).

Leave-one-team-out holds out one team $t$ at a time.
Test matches are those in which $t$ is home or away.
Training is the remaining matches, so the opponent's other fixtures stay in training.
Each match is scored once for each participating team that yields a valid fold, so a match typically contributes twice to the mean.
Uncertainty for per-KPI gaps and skill uses a team multinomial bootstrap: teams are redrawn with replacement, training weights are products of the two clubs' multiplicities, and out-of-bag matches form the test set.
Sport-level intervals for median gaps resample KPIs within sport.
The main text does not attach a $p$-value to one KPI's likelihood-ratio statistic.
Landmarks $\PEFLRaic$ and $\PEFLRlandmark$ on \cref{fig:e6_alignment} are references for feature choice and for a $\chi^2_1$ $5\%$ point, not discovery counts.

\clearpage

% ---------------------------------------------------------------------------
\refstepcounter{sisection}
\section*{SI Section \thesisection: Alignment and the Cost of a Redundant Coefficient}
\label{sec:si_alignment}

\noindent This section derives the two results used in the main text: the alignment condition behind the null $\beta_{\mathrm{A}}+\beta_{\mathrm{B}}=0$, and the finite-sample cost $1/(2n)$ of fitting the pair when the difference suffices.
Both are classical; we state them in the notation of this paper.

\subsection*{Alignment as a linear restriction}

Let $Y\in\{0,1\}$ indicate a home win, and let $A$ and $B$ be the home and away values of one KPI.
The pair model is
\begin{equation}
  \operatorname{logit}\Pr(Y=1\mid A,B)=\beta_{0}+\beta_{\mathrm{A}}A+\beta_{\mathrm{B}}B .
  \label{eq:si_pair}
\end{equation}
Writing $D=A-B$ and $S=A+B$, so that $A=(S+D)/2$ and $B=(S-D)/2$, the linear score becomes
\begin{equation}
  \beta_{\mathrm{A}}A+\beta_{\mathrm{B}}B
  =\tfrac12(\beta_{\mathrm{A}}-\beta_{\mathrm{B}})\,D
  +\tfrac12(\beta_{\mathrm{A}}+\beta_{\mathrm{B}})\,S .
  \label{eq:si_rotation}
\end{equation}
The difference model $\operatorname{logit}\Pr(Y=1\mid D)=\gamma_{0}+\gamma D$ is therefore the pair model with the single linear restriction $\beta_{\mathrm{A}}+\beta_{\mathrm{B}}=0$, with $\gamma=\beta_{\mathrm{A}}=-\beta_{\mathrm{B}}$.
When the restriction holds, $\Pr(Y=1\mid A,B)$ depends on $(A,B)$ only through $D$; within the model, $D$ is sufficient for the pair.
We call the KPI aligned in this case.
If the logistic model is misspecified, the same statements apply to the pseudo-true parameters that minimise expected log-loss \parencite{white1982}.

The restriction is invariant to centring and to rescaling $A$ and $B$ by a common factor, which is why both series are scaled by their pooled standard deviation.
Rescaling them by separate factors would change the null hypothesis.

Because the two models are nested and differ by one restriction, the likelihood-ratio statistic $\Lambda=2(\ell_{(\mathrm{A},\mathrm{B})}-\ell_{D})$ is asymptotically $\chi^{2}_{1}$ under the null for independent observations \parencite{wilks1938}.
Under local misalignment it is approximately noncentral $\chi^{2}_{1}$ with noncentrality $\lambda$, so that $\mathbb{E}[\Lambda]\approx 1+\lambda$.
Team recurrence makes matches dependent, so the $\chi^{2}_{1}$ reference is a landmark here rather than a calibrated null.

\subsection*{Expected excess test log-loss of a fitted model}

Let $\ell(\theta;x,y)$ be the log-likelihood of one observation under a regular parametric model with $k$ free parameters, and let $R(\theta)=-\mathbb{E}\,\ell(\theta;X,Y)$ be its expected log-loss on a new observation.
Let $\theta^{*}$ minimise $R$, and let $\hat\theta$ be the maximum-likelihood estimate from $n$ independent training observations.
A second-order expansion of $R$ about $\theta^{*}$ gives
\begin{equation}
  R(\hat\theta)\approx R(\theta^{*})
  +\tfrac12(\hat\theta-\theta^{*})^{\top}J\,(\hat\theta-\theta^{*}),
  \qquad J=\mathbb{E}\bigl[-\nabla^{2}\ell(\theta^{*};X,Y)\bigr].
\end{equation}
Since $\sqrt{n}(\hat\theta-\theta^{*})\to\mathcal{N}(0,J^{-1}KJ^{-1})$ with $K=\mathbb{E}\bigl[\nabla\ell\,\nabla\ell^{\top}\bigr]$ \parencite{white1982},
\begin{equation}
  \mathbb{E}\bigl[R(\hat\theta)\bigr]-R(\theta^{*})
  \approx\frac{\operatorname{tr}(J^{-1}K)}{2n}.
  \label{eq:si_excess}
\end{equation}
When the model is correctly specified, $K=J$ and the excess is $k/(2n)$.
This is the per-parameter penalty underlying Akaike's information criterion \parencite{akaike1973}; \cref{eq:si_excess} is the general form behind Takeuchi's criterion \parencite{konishi2008}.
No distributional assumption on the features is required beyond regularity.

\subsection*{The cost of fitting the pair when the difference suffices}

Suppose the KPI is aligned.
The pseudo-true pair parameter then satisfies the restriction, so the pair model (three parameters) and the difference model (two parameters) attain the same population risk $R(\theta^{*})$.
Applying \cref{eq:si_excess} to each model with correct specification,
\begin{equation}
  \mathbb{E}\bigl[\mathrm{LL}(D)-\mathrm{LL}(A,B)\bigr]
  \approx\frac{2}{2n}-\frac{3}{2n}=-\frac{1}{2n}.
  \label{eq:si_cost}
\end{equation}
Fitting the pair spends one redundant coefficient, and the difference is expected to have lower test log-loss by about $1/(2n)$ nats per match.
Under cross-validation, $n$ is the size of each training set, not the full sample.

If the KPI is not aligned, let $\Delta^{*}>0$ be the population log-loss lost by imposing the restriction.
The expected gap becomes
\begin{equation}
  \mathbb{E}\bigl[\mathrm{LL}(D)-\mathrm{LL}(A,B)\bigr]
  \approx\Delta^{*}-\frac{1}{2n},
  \label{eq:si_crossover}
\end{equation}
so the difference is still expected to predict better whenever the misalignment costs less than the redundant coefficient.
Under local alternatives the noncentrality of $\Lambda$ is $\lambda\approx2n\Delta^{*}$, so the crossover $\Delta^{*}=1/(2n)$ corresponds to $\lambda=1$, that is to $\mathbb{E}[\Lambda]\approx 2$.
This is Akaike's threshold for choosing between nested models that differ by one parameter, and it matches the asymptotic equivalence between Akaike's criterion and cross-validation \parencite{stone1977}.
For out-of-sample feature choice, $\Lambda=2$ is therefore the natural landmark; the $\chi^{2}_{1}$ $5\%$ point $3.84$ answers a different question (rejecting exact alignment).

\subsection*{Scope}

The expansions are first-order in $1/n$ and assume independent observations and a regular model.
Team recurrence and leave-one-team-out folds make matches dependent and folds unequal, so \cref{eq:si_cost,eq:si_crossover} are approximations for these leagues rather than exact predictions.
Under misspecification, $1/(2n)$ in \cref{eq:si_cost} is replaced by the difference in the two models' values of $\operatorname{tr}(J^{-1}K)/(2n)$.
